\section{The System Under Test}
\label{sec:system}

\texttt{agentic-data-contracts} is an open-source Python library that loads a
data contract and exposes a warehouse to an agent through nine governed
tools. The library is the instrument, not the contribution, and the
experiment is designed so that its identity does not matter: the hollow arm
runs the identical library, tools and rules with the prose removed.

\subsection{The contract}
\label{sec:contract}

The contract is two YAML files, 60{,}181 characters together.
\texttt{contract.yml} declares the five tables the agent may read, the eleven
operations it may not issue (\texttt{DELETE}, \texttt{DROP},
\texttt{UPDATE}, \texttt{INSERT}, \texttt{CREATE}, \texttt{ALTER} and five
more), and eight domains, each with a one-line summary and a prose
description; the \texttt{fees} description restates the manual's account of
every column of the fee-rule table. \texttt{semantic.yml} declares 14
metrics, each with a description and an executable
\texttt{sql\_expression}, the five tables with a description per column, and
one relationship. Figure~\ref{fig:contract} shows one metric as it appears in
the contract and in its hollow counterpart. Seven metric descriptions carry a
sentence beginning \texttt{INTERPRETATION}, marking where the contract
commits to a reading the manual does not force; Section~\ref{sec:threats}
returns to them.

The same file supplies both the semantics the agent reads and the rules the
tool layer enforces. The agent reads the first through
\texttt{lookup\_domain} and \texttt{lookup\_metric}; the tool layer checks
every statement against the second before it reaches the database.

\begin{figure*}[t]
\begin{minipage}[t]{0.63\textwidth}
\begin{lstlisting}
- name: fee_rule_matches_transaction
  description: >
    True when fee rule `f` applies to transaction `p` of merchant `m`, on
    every field that a single transaction can decide. The manual's rule is
    that "if a field is set to null it means that it applies to all possible
    values of that field", so each clause below is satisfied either by the
    field being null or by the field agreeing with the transaction. [...]
    INTERPRETATION: the list-typed fields (account_type, aci,
    merchant_category_code) are never null in the annexed `fees` data --
    they express "applies to all values" as an empty list -- so each clause
    accepts an empty list as the wildcard the manual describes.
  sql_expression: >
    (f.card_scheme IS NULL OR f.card_scheme = p.card_scheme)
    AND (f.is_credit IS NULL OR f.is_credit = p.is_credit)
    AND (f.aci IS NULL OR len(f.aci) = 0 OR list_contains(f.aci, p.aci))
    AND (f.intracountry IS NULL
         OR (f.intracountry = 1) = (p.issuing_country = p.acquirer_country))
    [... three more clauses: account_type, merchant_category_code,
     capture_delay mapped from days to the manual's bands]
  source_model: main.fees
  domains: [fees]
\end{lstlisting}
\end{minipage}\hfill
\begin{minipage}[t]{0.35\textwidth}
\begin{lstlisting}
- name: fee_rule_matches_transaction
  description: No definition is
    available for this metric.
  sql_expression: ''
  source_model: main.fees
  domains:
  - fees
\end{lstlisting}
\vspace{2pt}
{\small The same metric in the hollow contract. Name, source table and
domain membership survive; the description and the SQL expression do not.
Domain prose and column descriptions are emptied the same way.}
\end{minipage}
\caption{One of the 14 metrics of the frozen contract (left, abridged) and
its hollow counterpart (right). The contract arm can fetch the left with one
\texttt{lookup\_metric} call; the hollow arm fetches the right.}
\label{fig:contract}
\end{figure*}

\subsection{Nine tools and two-layer validation}
\label{sec:tools}

{\sloppy The tools fall into three groups. Four discover structure:
\texttt{describe\_table}, \texttt{preview\_table},
\texttt{lookup\_relationships} and \texttt{list\_metrics}. Three deliver
semantics: \texttt{lookup\_domain} returns a domain's prose,
\texttt{lookup\_metric} a metric's description and SQL expression, and
\texttt{trace\_metric\_impacts} the metrics that depend on a table. Two
handle queries: \texttt{inspect\_query} validates a statement and returns
either a pass or a reason, and \texttt{run\_query} validates and executes
it. Validation is two-layer. Every statement is parsed with
sqlglot~\cite{sqlglot} and checked statically against the table allow-list
and the forbidden operations; where a database adapter is available, a dry
run follows. A rejected statement returns its reason to the model rather
than an error to the user. The tool descriptions themselves carry 3{,}042
characters of procedural text (\texttt{inspect\_query} tells the model to
call \texttt{lookup\_metric} first), and that text is part of what the
hollow arm holds fixed.\par}

\subsection{The hollow contract}
\label{sec:hollow}

The hollow contract is generated from the frozen contract by a script, not
written by hand, so the only difference between the two artifacts is the one
the generator makes, and it is auditable as a diff. Emptied: domain summaries
and descriptions, metric descriptions, every \texttt{sql\_expression}
(the fee formula lives there), and all column and relationship descriptions.
Retained byte-for-byte: every name, every type, the structure, the table
allow-list, the forbidden operations and the identical nine-tool surface. An
agent that calls \texttt{lookup\_domain} against it is told \emph{``No
documentation is available for this domain.''}

The control is verified in both directions. The test suite asserts that no
6-gram of \texttt{manual.md} survives into the hollow contract, and that the
real contract does share 6-grams with it; the second half is what makes the
first meaningful. The transcripts confirm that the retrieval loop ran: on
\mGLM{} the hollow arm called \texttt{lookup\_domain} on 272 of 401 tasks and
received the placeholder 282 times. Prompt length is deliberately not held
constant (1{,}984 characters against the contract arm's 2{,}475): padding the
hollow arm with text that reads as content but carries none would be a
different treatment, not a cleaner control. The variable under test is
knowledge, not tokens.

\subsection{Provenance and freeze}
\label{sec:freeze}

The most obvious way to fake this result is to write a contract that answers
the benchmark. The contract's header records that every fact in it came from
\texttt{manual.md} and the dataset README and nothing else, which a reader
can verify by reading the 60{,}181-character contract against the
22{,}127-character manual. The artifact is content-hashed
(\texttt{sha256:e438ecf7\ldots}), was frozen roughly 35 commits before the
first run, and every result row records its digest, so a post-hoc edit would
invalidate every row.\footnote{On three of the four models the hollow arm's
rows carry the real contract's hash rather than the hollow artifact's, a
stamping defect that changes no result (\appref{app:harness}).} What cannot be verified from outside is that
no benchmark question was read while the contract was authored. We claim
provenance; we cannot prove intent.
